\section{时空优化}

\subsection{读入优化}

\verb|read<T>()|：读取一个值。\verb|read<int>()|、\verb|read<long long>()| 走 \texttt{getchar} 快读，\verb|read<std::string>()| 读取空白分隔的串，其余类型回落 \texttt{cin}（快读与 \texttt{cin} 不要混用）。

\begin{lstlisting}
namespace QuickStream {
    template<typename T>
    T read() {
        T x;
        std::cin >> x;
        return x;
    }

    template<typename T>
    inline T read_int() {
        T x = 0;
        bool f = 0;
        int ch = getchar();
        while (ch != EOF && (ch < '0' || ch > '9')) {
            if (ch == '-') f = 1;
            ch = getchar();
        }
        for (; ch >= '0' && ch <= '9'; ch = getchar())
            x = x * 10 + (ch - '0');
        return f ? -x : x;
    }

    template<>
    int read<int>() {
        return read_int<int>();
    }

    template<>
    long long read<long long>() {
        return read_int<long long>();
    }

    template<>
    std::string read<std::string>() {
        std::string s;
        int ch = getchar();
        while (ch != EOF && (ch == ' ' || ch == '\n' || ch == '\r' || ch == '\t')) {
            ch = getchar();
        }
        while (ch != EOF && ch != ' ' && ch != '\n' && ch != '\r' && ch != '\t') {
            s += (char)ch;
            ch = getchar();
        }
        return s;
    }
}
using QuickStream::read;
\end{lstlisting}

\subsection{指令集优化}

启用 \texttt{O3}、循环展开与常见位运算指令集，对大常数运算有明显加速；要求评测机为 \texttt{GCC} 且支持相应指令（不支持时 \texttt{pragma} 会被忽略），部分 OJ 禁用。

\begin{lstlisting}
#pragma GCC optimize("O3,unroll-loops,inline")
#pragma GCC target("avx2,bmi,bmi2,lzcnt,popcnt")
\end{lstlisting}
